\section*{Findings}
Offices with large and small 1944 staff stocks drew transferred workers from statistically similar organizational sources in both eras. Neither the share arriving from another ka in the same kyoku nor the share arriving from another kyoku differs significantly by stock group. The continuous-stock specifications reach the same conclusion and show no significant wartime change in either source-composition gradient.

The difference lies in transfer incidence rather than source composition. Larger offices had lower linked-worker transfer rates in peacetime. This negative stock gradient became substantially weaker during wartime. Accordingly, 1944 size predicts how frequently linked workers changed offices more clearly than it predicts whether transfers came from within or outside the destination kyoku.

\paragraph{Caveat.} The 1943 Shi--Fu merger changed many office identifiers. Only 11 drafted 1944 destination offices can be linked by \texttt{office\_id} to both eras. The estimates therefore describe organizationally stable offices and have limited statistical power. Extending the analysis requires an explicit predecessor-to-1944 office crosswalk.
